\documentclass[11pt]{article}
\usepackage[a4paper,margin=1in]{geometry}
\usepackage{amsmath,amsfonts,amssymb}
\usepackage{graphicx}
\usepackage{booktabs}
\usepackage{enumitem}
\usepackage[hidelinks]{hyperref}
\usepackage{microtype}

\title{Lazy ring + one directed shortcut: $\beta$ sweep, tail mechanism, and jump-over extension report}
\author{}
\date{}

\setlist[itemize]{itemsep=0pt,topsep=4pt,parsep=0pt,partopsep=0pt}
\setlength{\textfloatsep}{8pt plus 2pt minus 2pt}
\setlength{\floatsep}{6pt plus 2pt minus 2pt}
\setlength{\intextsep}{6pt plus 2pt minus 2pt}
\setlength{\abovecaptionskip}{4pt}
\setlength{\belowcaptionskip}{0pt}

\begin{document}
\maketitle

\section*{Key conclusions and key figures}
With $N=100,q=2/3,\rho=1$ and only $\beta$ varied, both peaks move to earlier times and the tail rate $\gamma$ increases with $\beta$;
at the same time, $K=4$ sustains a wider bimodal window and a deeper valley.
Figure~\ref{fig:beta-sweep-key-en} summarizes the peak and tail trends.

\begin{figure}[!htbp]
  \centering
  \begin{minipage}{0.48\linewidth}
    \centering
    \textbf{(a) Peak heights}\par
    \includegraphics[width=\linewidth]{outputs/beta_sweep_N100/peak_heights_vs_beta_N100.png}
  \end{minipage}\hfill
  \begin{minipage}{0.48\linewidth}
    \centering
    \textbf{(b) Peak times}\par
    \includegraphics[width=\linewidth]{outputs/beta_sweep_N100/peak_times_vs_beta_N100.png}
  \end{minipage}
  \vspace{4pt}
  \begin{minipage}{0.48\linewidth}
    \centering
    \textbf{(c) Tail $\gamma$}\par
    \includegraphics[width=\linewidth]{outputs/beta_sweep_N100/tail_gamma_vs_beta_N100.png}
  \end{minipage}\hfill
  \begin{minipage}{0.48\linewidth}
    \centering
    \textbf{(d) AW overlay for all $\beta$}\par
    \includegraphics[width=\linewidth]{outputs/beta_sweep_N100/aw_all_betas_N100.png}
  \end{minipage}
  \caption{Exact AW results under a $\beta$-only sweep: peak heights, peak times, tail-decay rate $\gamma$, and the full-$\beta$ overlay of $f(t)$; $\beta=0$ is the no-shortcut baseline and is shown as a standalone marker.}
  \label{fig:beta-sweep-key-en}
\end{figure}

At $\beta=0.01$, the representative $f(t)$ highlights the valley contrast:
$\texttt{hv\_over\_max}\approx0.34$ for $K=2$ versus about $0.15$ for $K=4$, indicating a deeper valley and clearer separation for $K=4$;
at the same $\beta$, $K=4$ has earlier $(t_1,t_2)$ and larger $\gamma$, consistent with faster time scales and lighter tails.
Figure~\ref{fig:beta-example-ft-en} shows the peak--valley window directly.

\begin{figure}[!htbp]
  \centering
  \begin{minipage}{0.48\linewidth}
    \centering
    \textbf{(a) $K=2$}\par
    \includegraphics[width=\linewidth]{outputs/mc_beta_sweep_N100/cases/N100_K2_beta0.01/N100_K2_beta0.f_t.pdf}
  \end{minipage}\hfill
  \begin{minipage}{0.48\linewidth}
    \centering
    \textbf{(b) $K=4$}\par
    \includegraphics[width=\linewidth]{outputs/mc_beta_sweep_N100/cases/N100_K4_beta0.01/N100_K4_beta0.f_t.pdf}
  \end{minipage}
  \caption{Representative exact $f(t)$ (AW), zoomed to the peak/valley window.}
  \label{fig:beta-example-ft-en}
\end{figure}

The tail mechanism is governed by the transient spectral radius, $\gamma=-\log(\lambda_{\max})$,
and $\log S(t)$ is close to linear in the tail (Fig.~\ref{fig:tail-logS-en}).

\begin{figure}[!htbp]
  \centering
  \begin{minipage}{0.48\linewidth}
    \centering
    \textbf{(a) $K=2$}\par
    \includegraphics[width=\linewidth]{outputs/tail_diag/log_survival_N100_K2_beta0.01.png}
  \end{minipage}\hfill
  \begin{minipage}{0.48\linewidth}
    \centering
    \textbf{(b) $K=4$}\par
    \includegraphics[width=\linewidth]{outputs/tail_diag/log_survival_N100_K4_beta0.01.png}
  \end{minipage}
  \caption{Log survival $\log S(t)$ at $\beta=0.01$, consistent with the spectral-radius tail rate $\gamma$.}
  \label{fig:tail-logS-en}
\end{figure}

MC classification confirms that jump-over appears only for $K=4$.
We sort trajectories by whether the shortcut is used at least once ($C0$ vs $C1p$) and whether a jump-over occurs ($J0$ vs $J1p$),
giving the four classes $\texttt{C0J0}$, $\texttt{C1pJ0}$, $\texttt{C0J1p}$, and $\texttt{C1pJ1p}$.
Under this lens, $K=2$ keeps $J$-related classes at zero across windows:
peak1 is essentially all $\texttt{C1pJ0}$, the valley is mostly $\texttt{C1pJ0}$ with a small $\texttt{C0J0}$ tail,
and peak2 is dominated by $\texttt{C0J0}$ with a smaller $\texttt{C1pJ0}$ share.
For $K=4$, peak1 splits into roughly two-thirds $\texttt{C1pJ0}$ and one-third $\texttt{C1pJ1p}$;
the valley shows a similar one-third $J1p$ share, while peak2 flips to $\texttt{C0}$-dominated with about a quarter $\texttt{C0J1p}$.
These $J1p$ paths correspond to delayed crossings without immediate absorption, and the curves vary only mildly with $\beta$ (Fig.~\ref{fig:mc-beta-classes-en}).

\begin{figure}[!htbp]
  \centering
  \begin{minipage}{0.32\linewidth}
    \centering
    \textbf{(a) peak1}\par
    \includegraphics[width=\linewidth]{outputs/mc_beta_sweep_N100/classes_peak1_vs_beta.png}
  \end{minipage}\hfill
  \begin{minipage}{0.32\linewidth}
    \centering
    \textbf{(b) valley}\par
    \includegraphics[width=\linewidth]{outputs/mc_beta_sweep_N100/classes_valley_vs_beta.png}
  \end{minipage}\hfill
  \begin{minipage}{0.32\linewidth}
    \centering
    \textbf{(c) peak2}\par
    \includegraphics[width=\linewidth]{outputs/mc_beta_sweep_N100/classes_peak2_vs_beta.png}
  \end{minipage}
  \caption{Conditional class probabilities vs $\beta$ for each window ($N=100$).}
  \label{fig:mc-beta-classes-en}
\end{figure}

To choose a representative bimodal setting, we apply $t_2/t_1\ge10$, $\texttt{hv\_over\_max}\le0.6$,
and $\texttt{h2\_over\_h1}\ge0.05$, selecting the smallest $\beta$ that satisfies both $K=2$ and $K=4$; the scan yields $\beta_\ast=0.01$.

Fixing $\beta_\ast=0.01$ and sweeping $N$ shows $t_2$ shifting right, $h_2$ decreasing, and $\gamma$ decreasing (heavier tails),
while $t_1$ stays roughly stable; see Figures~\ref{fig:n-sweep-peaks-en} and~\ref{fig:n-sweep-gamma-en}.

\begin{figure}[!htbp]
  \centering
  \begin{minipage}{0.48\linewidth}
    \centering
    \textbf{(a) $t_1$}\par
    \includegraphics[width=\linewidth]{outputs/N_sweep_beta_selected/t1_vs_N_beta0.01.png}
  \end{minipage}\hfill
  \begin{minipage}{0.48\linewidth}
    \centering
    \textbf{(b) $t_2$}\par
    \includegraphics[width=\linewidth]{outputs/N_sweep_beta_selected/t2_vs_N_beta0.01.png}
  \end{minipage}
  \vspace{4pt}
  \begin{minipage}{0.48\linewidth}
    \centering
    \textbf{(c) $h_1$}\par
    \includegraphics[width=\linewidth]{outputs/N_sweep_beta_selected/h1_vs_N_beta0.01.png}
  \end{minipage}\hfill
  \begin{minipage}{0.48\linewidth}
    \centering
    \textbf{(d) $h_2$}\par
    \includegraphics[width=\linewidth]{outputs/N_sweep_beta_selected/h2_vs_N_beta0.01.png}
  \end{minipage}
  \caption{Peak locations and heights vs $N$ at $\beta=0.01$.}
  \label{fig:n-sweep-peaks-en}
\end{figure}

\begin{figure}[!htbp]
  \centering
  \includegraphics[width=0.7\linewidth]{outputs/N_sweep_beta_selected/tail_gamma_vs_N_beta0.01.png}
  \caption{Tail decay rate $\gamma$ vs $N$ at $\beta=0.01$.}
  \label{fig:n-sweep-gamma-en}
\end{figure}

The MC sweep over $N$ sharpens the same contrast.
For $K=2$, everything remains in $J0$: peak1 is locked to $\texttt{C1pJ0}$, the valley is mostly $\texttt{C1pJ0}$,
and peak2 mixes $\texttt{C0J0}$ with $\texttt{C1pJ0}$ as $N$ changes.
For $K=4$, once the valley/second peak are well formed, peak1 and the valley hold a stable $\sim$one-third $J1p$ share
(largely $\texttt{C1pJ1p}$), while peak2 stays $\texttt{C0}$-dominated with about a quarter $\texttt{C0J1p}$,
showing that jump-over contributions persist across $N$ (Fig.~\ref{fig:mc-n-classes-en}).

\begin{figure}[!htbp]
  \centering
  \begin{minipage}{0.32\linewidth}
    \centering
    \textbf{(a) peak1}\par
    \includegraphics[width=\linewidth]{outputs/mc_N_sweep_beta_selected/classes_peak1_vs_N.png}
  \end{minipage}\hfill
  \begin{minipage}{0.32\linewidth}
    \centering
    \textbf{(b) valley}\par
    \includegraphics[width=\linewidth]{outputs/mc_N_sweep_beta_selected/classes_valley_vs_N.png}
  \end{minipage}\hfill
  \begin{minipage}{0.32\linewidth}
    \centering
    \textbf{(c) peak2}\par
    \includegraphics[width=\linewidth]{outputs/mc_N_sweep_beta_selected/classes_peak2_vs_N.png}
  \end{minipage}
  \caption{Class proportions vs $N$ at fixed $\beta=0.01$.}
  \label{fig:mc-n-classes-en}
\end{figure}

This extension builds on the main jump-over report (lazy\_jump) and uses the same model and notation.
The remainder gives the theory and valley definition, then the $\beta$ sweep, MC classes, $N$ sweep, MC-$N$, and reproducibility.

\section{Theory: generating functions, AW inversion, and tail decay}

This section provides the theoretical backbone for our exact pmf computation and tail diagnostics.
We consider a lazy random walk on an $N$-node ring with $K=2k$ neighbors: with probability $q$ the walker moves on the ring, choosing uniformly from steps $\pm 1,\dots,\pm k$, and with probability $1-q$ it stays put. A directed shortcut from $\mathrm{src}$ to $\mathrm{dst}$ is added.

\subsection{Spectral decomposition and the resolvent without a defect}
Let $P_0$ denote the defect-free transition matrix. For non-source nodes, the one-step transition is
\[
P_0(i\to i)=1-q,\qquad
P_0(i\to i\pm r)=\frac{q}{K},\quad r=1,\dots,k.
\]
Because $P_0$ is translation-invariant (a circulant matrix), its eigenvectors are Fourier modes
\[
\phi_\ell(x)=\exp\Big(\frac{2\pi i\ell x}{N}\Big),\qquad \ell=0,\dots,N-1,
\]
with eigenvalues
\begin{equation}
\lambda_\ell
=(1-q)+\frac{q}{K}\sum_{r=1}^{k}\Big(e^{2\pi i\ell r/N}+e^{-2\pi i\ell r/N}\Big)
=(1-q)+\frac{2q}{K}\sum_{r=1}^{k}\cos\Big(\frac{2\pi \ell r}{N}\Big).
\end{equation}

Define the resolvent (matrix generating function)
\[
S^{(0)}(z)=(I-zP_0)^{-1},\qquad |z|<1.
\]
By spectral decomposition,
\begin{equation}
S^{(0)}(z)_{x,y}
=\frac{1}{N}\sum_{\ell=0}^{N-1}
\frac{\exp\Big(\frac{2\pi i\ell(x-y)}{N}\Big)}{1-z\lambda_\ell}.
\end{equation}
One may further define a displacement form $Q(d,z):=S^{(0)}(z)_{y+d,y}$ to efficiently assemble defect quantities.

\subsection{Directed shortcut under the self-loop rule: a rank-one update}
In the \texttt{lazy\_selfloop} mode, the shortcut probability is
\[
p=\beta(1-q),
\]
and the same amount is removed from the self-loop probability at the source: $P(\mathrm{src}\to\mathrm{src})=(1-q)-p$ and $P(\mathrm{src}\to\mathrm{dst})=p$.
Hence the transition matrix is a rank-one perturbation of $P_0$:
\begin{equation}
P = P_0 + p\, e_{\mathrm{src}}(e_{\mathrm{dst}}-e_{\mathrm{src}})^\top,
\end{equation}
where $e_i$ is the standard basis vector. Therefore
\[
I-zP = (I-zP_0) - z p\, e_{\mathrm{src}}(e_{\mathrm{dst}}-e_{\mathrm{src}})^\top.
\]
Applying the Sherman--Morrison formula
\[
(A+uv^\top)^{-1}=A^{-1}-\frac{A^{-1}uv^\top A^{-1}}{1+v^\top A^{-1}u},
\]
one can obtain the required entries of the defect resolvent $S(z)$ from $S^{(0)}(z)$ with $O(1)$ additional scalar operations (as implemented in the repository pipeline).

\subsection{First absorption generating function under geometric absorption (\texorpdfstring{$\rho$}{rho})}
Let $\mathrm{tar}$ be the target. Upon each visit to $\mathrm{tar}$, absorption occurs with probability $\rho$, otherwise the process continues with probability $1-\rho$.
Let
\[
S_{i,j}(z)=\sum_{t\ge 0}\mathbb{P}(X_t=i\mid X_0=j)z^t
\]
be the visit generating function of the non-absorbed walk. Then the generating function of the first absorption pmf $f(t)=\mathbb{P}(T=t)$ for $t\ge 1$ is
\begin{equation}
\tilde f(z)=\sum_{t\ge 1}f(t)z^t
=\frac{\rho\,S_{\mathrm{tar},n_0}(z)}{(1-\rho)+\rho\,S_{\mathrm{tar},\mathrm{tar}}(z)}.
\end{equation}
For $\rho=1$, this reduces to
\[
\tilde f(z)=\frac{S_{\mathrm{tar},n_0}(z)}{S_{\mathrm{tar},\mathrm{tar}}(z)}.
\]

\subsection{Abate--Whitt (AW) inversion via FFT}
Coefficient extraction can be discretized on a circle. With $z_m=re^{2\pi i m/L}$ ($r<1$),
\begin{equation}
f(t)\approx r^{-t}\frac{1}{L}\sum_{m=0}^{L-1}\tilde f(z_m)e^{-2\pi i mt/L}.
\end{equation}
The right-hand side is an inverse DFT, thus FFT yields all $f(t)$ for $t=1,\dots,T_{\max}$ in $O(L\log L)$ time. We follow the repository's AW parameter selection and numerical stabilization.

\subsection{Exponential tail decay and the spectral-radius rate \texorpdfstring{$\gamma$}{gamma}}
Let $Q$ be the transient submatrix of the absorption chain. The survival function is
\[
S(t)=\mathbb{P}(T>t)=\mathbf{1}^\top Q^t e_{n_0}.
\]
By Perron--Frobenius theory for nonnegative matrices, there exists $\lambda_{\max}\in(0,1)$ such that
\[
S(t)\sim c\,\lambda_{\max}^t\quad (t\to\infty).
\]
We define the asymptotic exponential tail decay rate
\begin{equation}
\gamma=-\log(\lambda_{\max}),
\end{equation}
which provides a stable and comparable tail metric across $\beta$, $N$, and $K$. Numerically, we construct $Q$ and compute its spectral radius; for $\rho=1$ we remove the target state, and for $\rho<1$ we multiply the target row by $(1-\rho)$ to represent absorption leakage.

\section{Comparable definition of the valley and windowing}

Let $f(t)=\mathbb{P}(T=t)$ for $t\in\mathbb{N}^+$ denote the pmf of the first absorption time.
We detect candidate peaks as strict local maxima:
\begin{equation}
f(t)>f(t-1),\qquad f(t)>f(t+1),\qquad f(t)\ge h_{\min}.
\end{equation}
We further filter secondary peaks by a relative-height threshold:
$f(t)\ge \texttt{second\_rel\_height}\cdot \max_{u} f(u)$.

\subsection{V0: Default valley (minimum between the two peaks)}
Given the first two peaks in chronological order $t_1<t_2$, define
\begin{equation}
t_v=\arg\min_{t_1<t<t_2} f(t).
\end{equation}
This definition is parameter-light and guarantees that $t_v$ lies strictly between the two peaks,
facilitating comparisons across $\beta$ and across $K$.

\subsection{V1: Relative-distance window scale (for trajectory-type comparisons)}
Let $\Delta T=t_2-t_1$ be the peak separation and define the half window width
\begin{equation}
\Delta=\max\left(1,\left\lfloor \delta_{\mathrm{frac}}\Delta T\right\rfloor\right).
\end{equation}
Each sample is assigned to the nearest center in $\{t_1,t_v,t_2\}$ if
$\min_{c\in\{t_1,t_v,t_2\}} |T-c|\le \Delta$, otherwise it is assigned to \texttt{other}.
Because $\Delta$ scales proportionally with $\Delta T$, the windowing remains comparable across $\beta$.

\subsection{V2: Robust valley (excluding regions close to peaks)}
To avoid valleys too close to a peak (which may happen when peaks are close),
we may define
\begin{equation}
t_v=\arg\min_{t_1+\alpha\Delta T\le t\le t_2-\alpha\Delta T} f(t),
\end{equation}
with $\alpha\in(0,1/2)$ (e.g., $\alpha=0.1$). We use V2 as a sensitivity-analysis convention.

\section{Beta-only sweep at fixed \texorpdfstring{$N$}{N} (\texorpdfstring{$N=100$}{N=100})}
We fix $N=100$, $q=2/3$, $\rho=1$, and the \texttt{lazy\_selfloop} convention, and vary only $\beta$.
Figure~\ref{fig:beta-sweep-key-en} summarizes exact AW peak heights/locations and the tail-decay metric, together with an all-$\beta$ overlay of $f(t)$.

\paragraph{Peak--valley quantification}
We record the valley depth $\texttt{hv\_over\_max}=h_v/\max(h_1,h_2)$ and $\texttt{h2\_over\_h1}$ (see CSV).
Comparing $K=2$ and $K=4$ highlights the depth difference in the valley.
At $\beta=0.01$, $K=2$ has $\texttt{hv\_over\_max}\approx0.34$, while $K=4$ yields $\approx0.15$, indicating a deeper valley and clearer separation for $K=4$.
At the same $\beta$, $K=4$ shows earlier $t_1,t_2$ and larger $\gamma$, consistent with faster time scales and lighter tails.
Note that $K=2$ has $J\equiv0$ (no jump-over), while $K=4$ allows $J>0$, which will reshape class compositions in the MC sweep below.

\paragraph{Representative \texorpdfstring{$f(t)$}{f(t)} with improved readability}
Figure~\ref{fig:beta-example-ft-en} shows $\beta=0.01$ with the time axis focused on the peak--valley interval (avoiding the long vertical stretch near $t\approx0$), making the peak/valley contrast visually clear.

\paragraph{Tail decay diagnostics}
The asymptotic tail is controlled by $\gamma=-\log(\lambda_{\max})$.
Figure~\ref{fig:tail-logS-en} confirms the near-linear decay of $\log S(t)$.
At $\beta=0.01$, $K=2$ has $\gamma\approx3.89\times10^{-4}$ while $K=4$ has $\gamma\approx8.52\times10^{-4}$, matching the $\gamma$--vs--$\beta$ curves.

\subsection{Automatic choice of \texorpdfstring{$\beta$}{beta}}
We apply the same thresholds ($t_2/t_1\ge10$, $\texttt{hv\_over\_max}\le0.6$, $\texttt{h2\_over\_h1}\ge0.05$) and choose the smallest $\beta$ that satisfies them for both $K=2$ and $K=4$.
The scan yields $\beta_\ast=0.01$ (see \texttt{outputs/selected\_beta.txt}).

These exact AW results set the stage for the MC class analysis in the next section.

\section{MC sweep: four classes vs \texorpdfstring{$\beta$}{beta} (\texorpdfstring{$N=100$}{N=100})}
We use MC to estimate class probabilities within each window, defined by $(t_1,t_v,t_2)$ and $\Delta=\lfloor\delta_{\mathrm{frac}}(t_2-t_1)\rfloor$.
The four mutually exclusive classes are:
$\texttt{class\_C0J0}$ (no shortcut, no jump-over),
$\texttt{class\_C1pJ0}$ (shortcut, no jump-over),
$\texttt{class\_C0J1p}$ (no shortcut, jump-over),
$\texttt{class\_C1pJ1p}$ (shortcut and jump-over).
For $K=2$ we have $J\equiv0$, so $J$-related classes should be 0 or nearly 0.
Figure~\ref{fig:mc-beta-classes-en} reports the conditional class probabilities across windows as $\beta$ varies.

\paragraph{Interpretation highlights}
For $K=2$, $J$-related classes vanish, confirming that jump-over is absent.
For $K=4$, the valley and second-peak windows show larger shares of $J$-related classes, consistent with trajectories that cross the target without absorbing and hence contribute to delayed arrivals.

\begin{figure}[!htbp]
  \centering
  \begin{minipage}{0.32\linewidth}
    \centering
    \textbf{(a) peak1}\par
    \includegraphics[width=\linewidth]{outputs/mc_beta_sweep_N100/aux_peak1_vs_beta.png}
  \end{minipage}\hfill
  \begin{minipage}{0.32\linewidth}
    \centering
    \textbf{(b) valley}\par
    \includegraphics[width=\linewidth]{outputs/mc_beta_sweep_N100/aux_valley_vs_beta.png}
  \end{minipage}\hfill
  \begin{minipage}{0.32\linewidth}
    \centering
    \textbf{(c) peak2}\par
    \includegraphics[width=\linewidth]{outputs/mc_beta_sweep_N100/aux_peak2_vs_beta.png}
  \end{minipage}
  \caption{$\mathbb{P}(C\ge1\mid\text{window})$ and $\mathbb{P}(J\ge1\mid\text{window})$ vs $\beta$.}
  \label{fig:mc-beta-aux-en}
\end{figure}

With the $\beta$-direction class structure established, we now fix $\beta_\ast$ and examine how the slow time scale varies with $N$.

\section{Fix \texorpdfstring{$\beta=0.01$}{beta=0.01} and sweep \texorpdfstring{$N$}{N} (exact AW)}
Using the automatically selected $\beta_\ast=0.01$ and fixing $q=2/3$, $\rho=1$, and \texttt{lazy\_selfloop}, we sweep even $N$ and track peak locations/heights and tail decay.
Figure~\ref{fig:n-sweep-peaks-en} shows that $t_2$ shifts right as $N$ grows (a slower second time scale), while $h_2$ decreases; the tail rate $\gamma$ decreases with $N$, indicating a heavier tail, and the trend is summarized in Figure~\ref{fig:n-sweep-gamma-en}.

The next section complements this with MC class statistics across $N$.

\section{MC sweep: four classes vs \texorpdfstring{$N$}{N} (fixed \texorpdfstring{$\beta=0.01$}{beta=0.01})}
With $\beta=0.01$ fixed, we scan even $N$ and estimate class proportions within each window.
Recall: $C$ records whether the shortcut is used at least once ($C0$ vs $C1p$) and $J$ records whether a jump-over occurs ($J0$ vs $J1p$),
so the four classes are $\texttt{C0J0}$, $\texttt{C1pJ0}$, $\texttt{C0J1p}$, and $\texttt{C1pJ1p}$.
Figure~\ref{fig:mc-n-classes-en} summarizes the class proportions vs $N$.

\paragraph{Comparison highlights}
For $K=2$, $J$-related classes are essentially zero;
for $K=4$, $J$-related classes appear prominently in the valley and second-peak windows, reflecting jump-over contributions to delayed absorption.

\begin{figure}[!htbp]
  \centering
  \begin{minipage}{0.32\linewidth}
    \centering
    \textbf{(a) peak1}\par
    \includegraphics[width=\linewidth]{outputs/mc_N_sweep_beta_selected/aux_peak1_vs_N.png}
  \end{minipage}\hfill
  \begin{minipage}{0.32\linewidth}
    \centering
    \textbf{(b) valley}\par
    \includegraphics[width=\linewidth]{outputs/mc_N_sweep_beta_selected/aux_valley_vs_N.png}
  \end{minipage}\hfill
  \begin{minipage}{0.32\linewidth}
    \centering
    \textbf{(c) peak2}\par
    \includegraphics[width=\linewidth]{outputs/mc_N_sweep_beta_selected/aux_peak2_vs_N.png}
  \end{minipage}
  \caption{$\mathbb{P}(C\ge1\mid\text{window})$ and $\mathbb{P}(J\ge1\mid\text{window})$ vs $N$.}
  \label{fig:mc-n-aux-en}
\end{figure}

See the reproducibility section for the exact commands and parameters.

\section{Reproducibility (commands and outputs)}
Run the following inside \texttt{reports/lazy\_jump\_ext/}.

\paragraph{1. \texorpdfstring{$\beta$}{beta} sweep (exact AW)}
\begin{verbatim}
python code/beta_sweep_peaks_tail.py \
  --N 100 --Ks 2 4 --q 0.6666666667 --rho 1.0 --mode lazy_selfloop \
  --beta-min 0.0 --beta-max 0.2 --beta-num 21 \
  --outdir outputs/beta_sweep_N100
\end{verbatim}
Outputs: \texttt{outputs/beta\_sweep\_N100/beta\_sweep\_metrics\_N100.csv} and the three plots.

\paragraph{2. Automatic \texorpdfstring{$\beta$}{beta} selection}
\begin{verbatim}
python code/choose_beta.py \
  --csv outputs/beta_sweep_N100/beta_sweep_metrics_N100.csv \
  --Ks 2 4 --macro_ratio 10 --hv_over_max 0.6 --min_h2_over_h1 0.05 \
  --fallback_beta 0.02 \
  --out outputs/selected_beta.txt
\end{verbatim}

\paragraph{3. Fix \texorpdfstring{$\beta$}{beta} and sweep \texorpdfstring{$N$}{N} (exact AW)}
\begin{verbatim}
python code/N_sweep_peaks_tail.py \
  --K 2 4 --beta 0.01 --q 0.6666666667 --rho 1.0 --mode lazy_selfloop \
  --N-min 50 --N-max 300 --N-step 2 --only-even \
  --outdir outputs/N_sweep_beta_selected
\end{verbatim}

\paragraph{4. MC sweep over \texorpdfstring{$\beta$}{beta} (four classes)}
\begin{verbatim}
python code/mc_sweep_beta_classes.py \
  --N 100 --Ks 2 4 --q 0.6666666667 --rho 1.0 --mode lazy_selfloop \
  --betas 0.001 0.005 0.01 0.02 0.03 \
  --n-walkers 200000 --seed 0 \
  --outdir outputs/mc_beta_sweep_N100
\end{verbatim}

\paragraph{5. MC sweep over \texorpdfstring{$N$}{N} (fixed \texorpdfstring{$\beta$}{beta})}
\begin{verbatim}
python code/mc_sweep_N_classes.py \
  --beta 0.01 --Ks 2 4 --q 0.6666666667 --rho 1.0 --mode lazy_selfloop \
  --N-min 50 --N-max 300 --N-step 10 --only-even \
  --n-walkers 200000 --seed 0 \
  --outdir outputs/mc_N_sweep_beta_selected
\end{verbatim}

\paragraph{6. Tail diagnostics (\texorpdfstring{$\log S(t)$}{log S(t)})}
\begin{verbatim}
python code/tail_diagnostics.py \
  --N 100 --K 2 --beta 0.01 --q 0.6666666667 --rho 1.0 --mode lazy_selfloop \
  --max-steps-aw 12000 --outdir outputs/tail_diag
python code/tail_diagnostics.py \
  --N 100 --K 4 --beta 0.01 --q 0.6666666667 --rho 1.0 --mode lazy_selfloop \
  --max-steps-aw 12000 --outdir outputs/tail_diag
\end{verbatim}

\paragraph{7. Overlay all-\texorpdfstring{$\beta$}{beta} exact curves (optional)}
\begin{verbatim}
python code/plot_aw_all_betas.py \
  --N 100 --Ks 2 4 --q 0.6666666667 --rho 1.0 --mode lazy_selfloop \
  --beta-min 0.0 --beta-max 0.2 --beta-num 21 \
  --out outputs/beta_sweep_N100/aw_all_betas_N100.png
\end{verbatim}

\paragraph{8. Build the extension PDFs}
\begin{verbatim}
latexmk -xelatex -interaction=nonstopmode ring_lazy_jump_ext_cn.tex
latexmk -pdf -interaction=nonstopmode ring_lazy_jump_ext_en.tex
\end{verbatim}


\end{document}
